% BTS CIEL 1 — Services réseau — 2/4 : HTTP
% Compiler deux fois avec XeLaTeX depuis fichiers/.
\documentclass[aspectratio=169,11pt]{beamer}
\usepackage{fontspec}
\setmainfont{Latin Modern Sans}
\setsansfont{Latin Modern Sans}
\setmonofont{Latin Modern Mono}
\usepackage{polyglossia}
\setmainlanguage{french}
\usepackage{tikz}
\usetikzlibrary{arrows.meta,positioning,calc,shapes.geometric}
\usepackage{listings}
\usepackage{xcolor}
\definecolor{foret}{HTML}{1C322D}
\definecolor{creme}{HTML}{F8F3EE}
\definecolor{or}{HTML}{EBB552}
\definecolor{rose}{HTML}{F1CDBE}
\definecolor{vertclair}{HTML}{DDE9DF}
\definecolor{encre}{HTML}{24372E}
\setbeamercolor{normal text}{fg=creme,bg=foret}
\setbeamercolor{frametitle}{fg=or,bg=foret}
\setbeamercolor{title}{fg=creme}
\setbeamercolor{block title}{fg=foret,bg=or}
\setbeamercolor{block body}{fg=encre,bg=creme}
\setbeamercolor{alerted text}{fg=or}
\setbeamercolor{structure}{fg=or}
\setbeamertemplate{navigation symbols}{}
\setbeamertemplate{footline}{\hfill\color{or}\insertframenumber\hspace{0.5cm}\vspace{0.25cm}}
\setbeamertemplate{frametitle}{\vspace{0.35cm}\insertframetitle\par\vspace{0.15cm}\color{or}\rule{\linewidth}{0.6pt}}
\setlength{\parskip}{5pt}
\lstdefinestyle{webcode}{basicstyle=\ttfamily\small\color{encre},keywordstyle=\color{foret},stringstyle=\color{foret},
  backgroundcolor=\color{creme},frame=single,rulecolor=\color{or},framerule=1pt,
  xleftmargin=5pt,xrightmargin=5pt,columns=fullflexible,keepspaces=true,showstringspaces=false}
\newcommand{\question}[1]{\begin{beamercolorbox}[rounded=true,sep=12pt]{block body}\Large\textbf{#1}\end{beamercolorbox}}
\newcommand{\etiquette}[1]{{\color{or}\small\MakeUppercase{#1}}}
% Les schémas ont une surface claire, également lisible en noir et blanc.
\tikzset{
  boite/.style={draw=encre,fill=creme,text=encre,rounded corners,
    align=center,text width=2.8cm,minimum height=1cm},
  flux/.style={-{Stealth},thick,draw=encre},
  every node/.style={text=encre,font=\small}}
\newenvironment{schema}{%
  \begin{beamercolorbox}[sep=8pt,rounded=true]{block body}
  \centering\begin{tikzpicture}[>=Stealth]
}{\end{tikzpicture}\end{beamercolorbox}}
\title{Le protocole HTTP}
\subtitle{BTS CIEL 1 · Services réseau · 2/4}
\author{}
\date{}
\begin{document}

\begin{frame}[plain]
  \vfill
  \etiquette{Demander et recevoir une page Web}\par\bigskip
  {\Huge\bfseries Le protocole HTTP}\par\bigskip
  {\Large BTS CIEL 1 · Services réseau · 2/4}\par\bigskip
  \begin{schema}
    \node[boite] (c) {Navigateur\\demande une page};
    \node[boite,right=4cm of c] (s) {Serveur HTTP\\renvoie un contenu};
    \draw[flux] (c)--node[below=3mm]{HTTP} (s);
  \end{schema}
  \vfill
\end{frame}

\begin{frame}{Rappel : qui parle, avec quelles règles ?}
  Pour ouvrir une page, le navigateur contacte un logiciel serveur.
  \begin{beamercolorbox}[sep=10pt]{block body}
  \renewcommand{\arraystretch}{1.4}
  \begin{tabular}{ll}
    \textbf{Client} & Le navigateur demande une ressource.\\
    \textbf{Serveur} & Le logiciel HTTP répond à la demande.\\
    \textbf{Protocole} & HTTP définit les règles des messages.\\
    \textbf{Port} & Le numéro aide à retrouver l'application.\\
  \end{tabular}
  \end{beamercolorbox}
  \begin{block}{Méthode}
    1. Repérez qui demande.\quad 2. Repérez qui répond.\\
    3. Nommez le protocole, puis le port du service.
  \end{block}
\end{frame}

\begin{frame}{Une adresse saisie, une page affichée}
  Vous tapez une adresse dans le navigateur. Une page s'affiche.
  Comment le serveur sait-il quelle page vous voulez ?
  \begin{schema}
    \node[boite] (c) {Navigateur\\« Je veux l'accueil. »};
    \node[boite,right=4cm of c] (s) {Serveur Web\\« Voici la page. »};
    \draw[flux] (c)--node[below=3mm]{langage commun} (s);
  \end{schema}
  \begin{block}{Premier principe}
    Sans règles communes, la demande serait incompréhensible.\\
    \textbf{HTTP} organise les demandes et les réponses du Web.
  \end{block}
\end{frame}

\begin{frame}{Exemple : PC0 ouvre la page d'accueil}
  \etiquette{Plateforme : Cisco Packet Tracer}\par
  Dans cet exemple, PC0 ouvre \texttt{http://192.168.1.251}.
  \begin{schema}
    \node[boite,text width=2.5cm] (c) at (0,0)
      {PC0\\navigateur};
    \node[boite,text width=2cm] (sw) at (4.2,0) {Switch};
    \node[boite,text width=3cm] (s) at (8.7,0)
      {Server1\\HTTP\\192.168.1.251};
    \draw[thick] (c)--(sw);
    \draw[thick] (sw)--(s);
  \end{schema}
  \begin{block}{Méthode de lecture}
    1. PC0 demande la page d'accueil.\\
    2. Server1 renvoie le contenu de cette page.\\
    3. Le navigateur interprète le HTML et affiche la page.
  \end{block}
\end{frame}

\begin{frame}{Lire une URL : où et quoi demander ?}
  Une URL est l'adresse d'une ressource Web.
  \begin{schema}
    \node[boite,text width=2cm] at (0,0)
      {\texttt{http://}\\protocole};
    \node[boite,text width=4cm] at (3.8,0)
      {\texttt{192.168.1.251}\\hôte : serveur visé};
    \node[boite,text width=3cm] at (8.5,0)
      {\texttt{/index.html}\\chemin : ressource};
  \end{schema}
  \begin{block}{Méthode}
    1. Avant \texttt{://}, lisez le protocole : \texttt{http}.\\
    2. Jusqu'au prochain \texttt{/}, repérez l'hôte.\\
    3. Lisez le chemin : ici \texttt{/index.html}.
  \end{block}
  Avec \texttt{http}, sans port écrit, le port serveur est \textbf{80}.
\end{frame}

\begin{frame}{Avant HTTP : établir une connexion TCP}
  Comme au téléphone, on établit le contact avant de parler.
  Ici, nous étudions \textbf{HTTP/1.1 sur TCP}.
  \begin{schema}
    \node[align=center] at (0,3.2)
      {\textbf{Client}\\port éphémère 50532};
    \node[align=center] at (8,3.2)
      {\textbf{Serveur HTTP}\\port 80};
    \draw[dashed] (0,2.7)--(0,-.1);
    \draw[dashed] (8,2.7)--(8,-.1);
    \draw[flux] (0,2.3)--node[below]{1. SYN : demande de connexion} (8,2.3);
    \draw[flux] (8,1.3)--node[below]{2. SYN-ACK : accord du serveur} (0,1.3);
    \draw[flux] (0,.3)--node[below]{3. ACK : confirmation du client} (8,.3);
  \end{schema}
  \small Lisez de haut en bas, puis suivez les flèches.
  50532 est un exemple ; le port client peut changer.
\end{frame}

\begin{frame}[fragile]{La requête : demander une ressource}
  Une connexion existe. Le navigateur précise maintenant ce qu'il demande.
  \begin{lstlisting}[style=webcode,language={},showlines=true]
GET /index.html HTTP/1.1
Host: 192.168.1.251

\end{lstlisting}
  Un en-tête est une information qui accompagne le message.\par
  \begin{block}{Méthode : lire le texte HTTP/1.1}
    1. \texttt{GET} : demander ; \texttt{/index.html} : chemin demandé.\\
    2. \texttt{HTTP/1.1} : version utilisée pour cet échange.\\
    3. \texttt{Host} : hôte visé, ici \texttt{192.168.1.251}.\\
    4. La ligne vide termine les en-têtes.
  \end{block}
\end{frame}

\begin{frame}{GET ou POST : quelle demande ?}
  Vous consultez une page, puis vous envoyez un formulaire.
  Ce sont deux actions différentes.
  \begin{beamercolorbox}[sep=10pt]{block body}
  \renewcommand{\arraystretch}{1.7}
  \begin{tabular}{ll}
    \textbf{GET} & Demander une ressource, par exemple une page.\\
    \textbf{POST} & Envoyer des données, par exemple un formulaire.\\
  \end{tabular}
  \end{beamercolorbox}
  \begin{block}{Méthode}
    1. Identifiez l'action : consulter ou envoyer des données.\\
    2. Associez-la à la méthode : ici GET ou POST.
  \end{block}
  \small Chiffrer, c'est rendre les données illisibles pour qui n'a pas la clé.\par
  POST seul ne chiffre pas les données.
\end{frame}

\begin{frame}[fragile]{La réponse : résultat, puis contenu}
  Le navigateur doit connaître le résultat et la nature du contenu reçu.\par
  \textbf{Extrait pédagogique d'une réponse HTTP/1.1 :}
  \begin{lstlisting}[style=webcode,language={}]
HTTP/1.1 200 OK
Content-Type: text/html

<html><body><h1>Bienvenue</h1></body></html>
\end{lstlisting}
  \begin{block}{Méthode}
    1. Lisez la ligne d'état : \texttt{200 OK}, la demande a réussi.\\
    2. Lisez \texttt{Content-Type} : le contenu est du HTML.\\
    3. Après la ligne vide commence le corps : ici le code HTML.
  \end{block}
  \small Les autres en-têtes sont omis dans cet extrait.
\end{frame}

\begin{frame}{Relier la requête et la réponse}
  Le navigateur demande \texttt{/index.html}. Le serveur répond \texttt{200 OK}.
  \begin{schema}
    \node[boite] (c) {Client\\navigateur};
    \node[boite,right=5.4cm of c] (s) {Serveur HTTP\\fournit la ressource};
    \draw[flux] ([yshift=3mm]c.east)--
      node[above]{GET \texttt{/index.html}} ([yshift=3mm]s.west);
    \draw[flux,dashed] ([yshift=-3mm]s.west)--
      node[below]{200 OK + contenu HTML} ([yshift=-3mm]c.east);
  \end{schema}
  \begin{block}{Méthode}
    1. Repérez la ressource dans la requête.\\
    2. Repérez le résultat dans la réponse.\\
    3. Distinguez le message HTTP du contenu qu'il transporte.
  \end{block}
\end{frame}

\begin{frame}{Les codes d'état : quel résultat ?}
  Une page absente ne peut pas être renvoyée comme une page trouvée.
  Le code d'état indique le résultat de la demande.
  \begin{beamercolorbox}[sep=9pt]{block body}
  \renewcommand{\arraystretch}{1.5}
  \small
  \begin{tabular}{lll}
    \textbf{Famille} & \textbf{Sens} & \textbf{Exemple}\\\hline
    2xx & Succès & \texttt{200 OK}\\
    3xx & Redirection & \texttt{301 Moved Permanently}\\
    4xx & Erreur côté demande & \texttt{404 Not Found}\\
    5xx & Erreur côté serveur & \texttt{500 Internal Server Error}\\
  \end{tabular}
  \end{beamercolorbox}
  \small 301 : ressource déplacée définitivement.
  404 : ressource non trouvée. 500 : erreur interne du serveur.
\end{frame}

\begin{frame}{Lire un code sans tout mémoriser}
  Vous demandez \texttt{/photo.png}. Le serveur répond \texttt{404 Not Found}.
  \begin{block}{Méthode}
    1. Lisez le premier chiffre : \textbf{4}, famille des erreurs côté demande.\\
    2. Lisez le code complet : \textbf{404}, ressource non trouvée.\\
    3. Vérifiez le chemin demandé et l'existence de la ressource.
  \end{block}
\end{frame}

\begin{frame}{Une page peut nécessiter plusieurs requêtes}
  Le HTML référence ici un fichier CSS et une image à récupérer.
  \begin{schema}
    \node[boite,text width=3.4cm] at (0,1.3)
      {1. GET \texttt{/index.html}\\reçoit le HTML};
    \node[boite,text width=3.4cm] at (6.5,2.0)
      {2. GET \texttt{/style.css}\\reçoit la mise en forme};
    \node[boite,text width=3.4cm] at (6.5,.6)
      {3. GET \texttt{/logo.png}\\reçoit l'image};
    \draw[flux] (1.9,1.3)--(4.6,2.0);
    \draw[flux] (1.9,1.3)--(4.6,.6);
  \end{schema}
  \begin{block}{Méthode sur cet exemple}
    1. Comptez le document HTML : une requête.\\
    2. Ajoutez les deux fichiers référencés : \textbf{trois requêtes au total}.
  \end{block}
  \small CSS et image peuvent être demandés en parallèle.
\end{frame}

\begin{frame}{HTTP en clair : un formulaire peut être lu}
  Vous envoyez un mot de passe avec un formulaire en HTTP.
  Une personne capable d'intercepter le trafic peut lire les données.
  \begin{schema}
    \node[boite] (c) {Navigateur\\formulaire envoyé};
    \node[boite,right=4.8cm of c] (s) {Serveur HTTP\\formulaire reçu};
    \draw[flux] (c)--node[below=3mm]{mot de passe lisible} (s);
  \end{schema}
  \begin{block}{Méthode}
    1. Repérez \texttt{http://} dans l'adresse.\\
    2. Considérez les données HTTP comme non chiffrées.\\
    3. Pour envoyer des données sensibles, utilisez HTTPS.
  \end{block}
\end{frame}

\begin{frame}{HTTPS : protéger l'échange HTTP}
  \textbf{TLS} protège HTTP (\texttt{https://}) : il chiffre les données.\par
  Il vérifie qu'on parle au vrai serveur : c'est l'authentification du serveur.
  \begin{schema}
    \node[boite] (c) {Navigateur};
    \node[boite,right=4.8cm of c] (s) {Serveur HTTPS\\port usuel 443};
    \draw[flux] (c)--node[below=3mm]{échange HTTP chiffré} (s);
  \end{schema}
  \begin{block}{Méthode dans le navigateur}
    1. Vérifiez \texttt{https://} et les informations de sécurité.\\
    2. Le cadenas, selon le navigateur, indique une connexion protégée.\\
    3. Il ne garantit pas l'honnêteté du site ; vérifiez aussi son adresse.
  \end{block}
  \small Sans port écrit : HTTPS utilise \textbf{443}, HTTP utilise \textbf{80}.
\end{frame}

\begin{frame}{Côté serveur : fournir les fichiers Web}
  Un logiciel comme \textbf{Apache} ou \textbf{Nginx} peut servir les fichiers d'un dossier.
  Par exemple, \texttt{/index.html} correspond à une page de ce site.\par
  \etiquette{Plateforme : Cisco Packet Tracer, Server-PT}\par
  \begin{block}{Méthode dans le simulateur}
    1. Ouvrez le serveur, puis \textbf{Services > HTTP}.\\
    2. Repérez le service HTTP et son état \textbf{On / Off}.\\
    3. Utilisez le gestionnaire de fichiers (\textbf{File Manager})
    pour éditer une page, puis enregistrez-la.
  \end{block}
  \small Le simulateur propose son propre service HTTP.

\end{frame}

\begin{frame}{Séance suivante : retenir un nom plutôt qu'une IP}
  Vous devez mémoriser \texttt{192.168.1.251} pour retrouver le site.
  Avec plusieurs sites, ces nombres deviennent difficiles à retenir.
  \begin{schema}
    \node[boite,text width=4cm] (n) {Nom facile à retenir\\par exemple \texttt{www.site.com}};
    \node[boite,text width=4cm,right=3cm of n] (i)
      {Adresse du serveur\\192.168.1.251};
    \draw[flux,dashed] (n)--node[below=3mm]{quelle association ?} (i);
  \end{schema}
  \begin{block}{Question pour la prochaine séance}
    Comment associer un nom à une adresse IPv4 ? Nous étudierons \textbf{DNS}.
  \end{block}
  \small Nom utilisé comme illustration ; aucune configuration DNS n'est supposée ici.
\end{frame}

\begin{frame}{À vous : décomposer une URL}
  \question{Dans \texttt{http://192.168.2.10/notice.html}, repérez le protocole,
    l'hôte et le chemin. Quel est le port serveur implicite ?}
  \bigskip
  \etiquette{Exemple d'entraînement hors TP}\par
  \small On suppose ici un service Web sur cette machine.
\end{frame}

\begin{frame}{Correction : lire les trois parties}
  \begin{block}{Décomposition}
    Protocole : \texttt{http}.\\
    Hôte : \texttt{192.168.2.10}.\\
    Chemin : \texttt{/notice.html}.\\
    Port serveur implicite : \textbf{80}.
  \end{block}
  \begin{enumerate}
    \item Coupez à \texttt{://} pour lire le protocole.
    \item Repérez l'hôte avant le prochain \texttt{/}.
    \item Gardez le \texttt{/} au début du chemin.
  \end{enumerate}
\end{frame}

\begin{frame}{À vous : interpréter une réponse}
  \question{Vous demandez \texttt{/cours.html} et recevez \texttt{500 Internal Server Error}.
    Que signifie ce code ? Le serveur HTTP a-t-il répondu ?}
  \bigskip
  \etiquette{Appuyez-vous sur la famille, puis sur le code complet}
\end{frame}

\begin{frame}{Correction : une erreur interne du serveur}
  \begin{block}{Interprétation}
    \textbf{500} signifie qu'une erreur interne du serveur empêche de traiter la demande.\\
    \textbf{Oui}, le serveur HTTP a répondu : il a envoyé ce code.
  \end{block}
  \begin{enumerate}
    \item Repérez la famille \textbf{5xx} : erreur côté serveur.
    \item Précisez le résultat : erreur interne lors du traitement de la demande.
    \item Faites vérifier le fonctionnement du logiciel serveur.
  \end{enumerate}
\end{frame}

\begin{frame}{À vous : reconnaître une connexion protégée}
  \question{Une URL commence par \texttt{https://}, sans numéro de port.
    Quel port serveur est utilisé par défaut ? Que protège HTTPS ?}
  \bigskip
  \etiquette{Distinguez le port du service et la protection de l'échange}
\end{frame}

\begin{frame}{Correction : HTTPS et le port 443}
  \begin{block}{Réponse}
    Port serveur par défaut : \textbf{443}.\\
    HTTPS protège les échanges HTTP grâce à \textbf{TLS} :\\
    chiffrement des données et authentification du serveur.
  \end{block}
  \begin{enumerate}
    \item Lisez le protocole au début de l'URL : \texttt{https}.
    \item Sans port indiqué, retenez \textbf{443}.
    \item Vérifiez l'adresse du site et les informations de sécurité.
  \end{enumerate}
\end{frame}

\begin{frame}{À retenir}
  \begin{enumerate}
    \item HTTP organise les requêtes du client et les réponses du serveur.
    \item Une URL indique le protocole, l'hôte et le chemin demandé.
    \item HTTP/1.1 utilise TCP ; le port HTTP usuel est \textbf{80}.
    \item Une réponse indique un code d'état et peut contenir une ressource.
    \item HTTPS protège HTTP avec TLS ; son port usuel est \textbf{443}.
  \end{enumerate}
  \bigskip
  \etiquette{Séance suivante : DNS — associer un nom à une adresse}
\end{frame}

\begin{frame}{Sources : HTTP et Web}
  \small
  \setlength{\parskip}{2pt}
  \begin{itemize}
    \item \textbf{RFC 9110 — Sémantique HTTP}\\
      \texttt{https://www.rfc-editor.org/rfc/rfc9110.html}\\
      URL (§4), méthodes (§9), codes d'état (§15).
    \item \textbf{RFC 9112 — HTTP/1.1}\\
      \texttt{https://www.rfc-editor.org/rfc/rfc9112.html}
    \item \textbf{MDN — HTTP}\\
      \texttt{https://developer.mozilla.org/fr/docs/Web/HTTP}
    \item \textbf{RFC 9293 — TCP}, connexion (§3.5)\\
      \texttt{https://www.rfc-editor.org/rfc/rfc9293.html}
  \end{itemize}
  \begin{block}{Supports de l'enseignant}
    « Protocoles réseaux », partie HTTP ;\\
    « TP DHCP, HTTP, DNS, FTP », pages 2--3.
  \end{block}
\end{frame}

\end{document}
